%% ACM conference paper skeleton (acmart, sigconf format)
%% Advanced Telerobotics Research Lab, Department of Computer Science, Kent State University
%%
%% Build:  tectonic main.tex        (or: latexmk -pdf main.tex)
%% Guide:  ../../../research/writing-papers.md
%%
%% Before you submit:
%%   1. Check the venue's call for papers for the format (sigconf or the
%%      single-column manuscript format), page or word limits and anonymity rules.
%%   2. Replace every [bracketed placeholder]. Search the file for "[".
%%   3. Generate your own CCS concepts at https://dl.acm.org/ccs/ccs.cfm and paste
%%      the CCSXML block and \ccsdesc lines below.

%% ---------------------------------------------------------------------------
%% Choose ONE \documentclass line.
%% ---------------------------------------------------------------------------
%% 1. Double-blind submission: line numbers on, authors and acknowledgments hidden.
\documentclass[sigconf,review,anonymous]{acmart}
%% 2. Single-blind submission: line numbers on, authors shown.
% \documentclass[sigconf,review]{acmart}
%% 3. Camera-ready, after acceptance (add the rights commands ACM sends you).
% \documentclass[sigconf]{acmart}
%% 4. Venues that ask for the single-column review format.
% \documentclass[manuscript,review,anonymous]{acmart}

%% ---------------------------------------------------------------------------
%% Rights and conference metadata
%% ---------------------------------------------------------------------------
%% After acceptance, ACM's rights form gives you the exact commands for your
%% paper (copyright or license type, DOI, ISBN, conference line). Replace this
%% whole block with them. Until then these placeholders keep the layout right.
\setcopyright{acmlicensed}
\copyrightyear{[YYYY]}
\acmYear{[YYYY]}
\acmDOI{XXXXXXX.XXXXXXX}
\acmConference[{[VENUE 'YY]}]{[Full Conference Name]}{[Dates]}{[City, Country]}
\acmISBN{978-X-XXXX-XXXX-X/YYYY/MM}

\begin{document}

%% ---------------------------------------------------------------------------
%% Title and authors
%% ---------------------------------------------------------------------------
\title{[Paper Title in the Venue's Capitalization Style]}

\author{[First Author Name]}
\email{[username]@kent.edu}
\orcid{0000-0000-0000-0000} % [your ORCID iD]; ACM requires one for every author
\affiliation{%
  \department[0]{Advanced Telerobotics Research Lab}
  \department[1]{Department of Computer Science} % [1] = the lab sits inside the department
  \institution{Kent State University}
  \city{Kent}
  \state{OH}
  \country{USA}
}

\author{[Second Author Name]}
\email{[username]@kent.edu}
\affiliation{%
  \department[0]{Advanced Telerobotics Research Lab}
  \department[1]{Department of Computer Science}
  \institution{Kent State University}
  \city{Kent}
  \state{OH}
  \country{USA}
}

\author{[Faculty Advisor Name]}
\email{[username]@kent.edu}
\affiliation{%
  \department[0]{Advanced Telerobotics Research Lab}
  \department[1]{Department of Computer Science}
  \institution{Kent State University}
  \city{Kent}
  \state{OH}
  \country{USA}
}

%% Short author list for page headers.
\renewcommand{\shortauthors}{[Surname] et al.}

\begin{abstract}
[One paragraph, usually 150 to 250 words; check the call for papers. State the
problem, what you built or studied, how you evaluated it and the main result
with its conditions (for example, ``in simulation'' or ``on a Unitree Go2 in the
lab'').]
\end{abstract}

%% CCS concepts: replace with your own from https://dl.acm.org/ccs/ccs.cfm
\begin{CCSXML}
<ccs2012>
   <concept>
       <concept_id>10010520.10010553.10010554</concept_id>
       <concept_desc>Computer systems organization~Robotics</concept_desc>
       <concept_significance>500</concept_significance>
   </concept>
   <concept>
       <concept_id>10003120.10003121</concept_id>
       <concept_desc>Human-centered computing~Human computer interaction (HCI)</concept_desc>
       <concept_significance>300</concept_significance>
   </concept>
</ccs2012>
\end{CCSXML}

\ccsdesc[500]{Computer systems organization~Robotics}
\ccsdesc[300]{Human-centered computing~Human computer interaction (HCI)}

\keywords{[telerobotics], [human-robot interaction], [keyword three]}

\maketitle

%% ---------------------------------------------------------------------------
\section{Introduction}
%% ---------------------------------------------------------------------------
[Motivate the problem and state your contributions as a short list.]
Robot software in the lab builds on ROS~2~\cite{macenski2022ros2,ros2jazzydocs},
and many of our controllers are first tested in simulation~\cite{todorov2012mujoco}.

\section{Related Work}
[Group prior work by approach and say how yours differs.]

\section{Method}
[Describe the system. Define every symbol the first time it appears.]
\begin{equation}
  \mathbf{u}_t = K\,(\mathbf{x}^{*}_t - \mathbf{x}_t)
  \label{eq:control}
\end{equation}
Equation~\eqref{eq:control} is a placeholder.

\section{Study}
[Participants (with the IRB protocol number once approved), apparatus,
procedure, measures and analysis. Say whether each robot run was
teleoperated, autonomous or scripted.]

%% Sample figure. Replace the \fbox placeholder with
%%   \includegraphics[width=\linewidth]{figures/[file].pdf}
%% \Description is required: it is the figure's alt text for screen readers.
\begin{figure}[t]
  \centering
  \fbox{\parbox[c][1.4in][c]{0.92\linewidth}{\centering Figure placeholder}}
  \caption{[One-sentence caption that states what the figure shows and under
  which conditions, for example ``Teleoperated trials, $n=[N]$'']}
  \Description{[Alt text: describe what the figure shows for a reader who
  cannot see it, including the main trend.]}
  \label{fig:overview}
\end{figure}

Figure~\ref{fig:overview} shows [what the reader should notice].

\begin{table}[t]
  \caption{[Table caption above the table, as ACM style requires]}
  \label{tab:results}
  \begin{tabular}{lcc}
    \toprule
    Condition & [Metric 1] & [Metric 2] \\
    \midrule
    {[Baseline]} & {[0.00]} & {[0.00]} \\
    {[Ours]}     & {[0.00]} & {[0.00]} \\
    \bottomrule
  \end{tabular}
\end{table}

\section{Discussion and Limitations}
[What the results do and do not show.]

\section{Conclusion}
[Two or three sentences.]

%% The acks environment is hidden automatically in anonymous mode.
\begin{acks}
%% Name funders exactly as the award record does, and only after the PI
%% confirms which award funded this work.
This work was supported in part by [Sponsor] under [Grant No.].
%% Generative AI: ACM's authorship policy (updated May 14, 2026) asks you to
%% describe AI used to do the research (data, code, analysis, figures) in the
%% Methods section; disclosing AI used only to help write is no longer
%% required. The venue may ask for more (see research/standards.md).
\end{acks}

\bibliographystyle{ACM-Reference-Format}
\bibliography{references}

\end{document}
